% mage-auto-heading: problem
% mage-auto-heading: conceptual-design
% mage-auto-heading: engineering-models
% mage-auto-heading: realization
% mage-auto-heading: connect-the-models-to-their-consumers
% mage-auto-heading: select-models-for-fixed-tasks
% mage-auto-heading: exercise-an-alternative-engine
% mage-auto-heading: evaluate-the-result
% mage-auto-heading: adaptation-and-limits
% Evidence: book/data/local-evidence/blueprints-11202026/evidence-ledger.md; DocAble revision 563cdae2316df8798391d17c794ae7f662410321.
% Authoritative LaTeX manuscript; Markdown is a generated compatibility projection.

\section{Problem}\label{heading:problem}

A system that relies on AI reasoning need not commit its architecture to a particular foundation model or reasoning strategy. Models differ in cost, latency, reliability, and capability, and these properties change as the technology advances. An engineering design that embeds one model, prompt, or orchestration strategy throughout the implementation makes alternatives expensive to evaluate. A better arrangement separates the reasoning that proposes work from the machinery that carries out and evaluates its effects.

\section{Conceptual Design}\label{heading:conceptual-design}

% Evidence: BP-SUBSTITUTION-CONTRACTS | source: book/data/local-evidence/evidentiary-basis-11202026/integration/records/BP-SUBSTITUTION-CONTRACTS.md | kind: implementation-and-design
DocAble supplies examples for two distinct substitution decisions. Model selection chooses which foundation model performs a defined semantic task. Replacing the reasoning engine changes how the system decomposes and sequences work: a general-purpose agent can take over planning decisions made by a graph of specialized remediation passes. This blueprint separates both decisions from the domain contracts that define permitted changes and acceptable results.

\autoref{handbook-substitutable-reasoning} illustrates the proposed architecture. Each reasoning engine inspects the artifact and proposes supported changes through a typed operation interface. An application component resolves targets, checks applicable preconditions, performs mutations, and records their effects. Separate evaluators inspect the resulting artifact. Hold the permitted operations and acceptance requirements fixed when comparing engines, so differences in the results can be investigated as differences in reasoning rather than changes to the task.

\begin{figure}[!htb]
\centering
\includegraphics[width=\linewidth]{book/latex/assets/handbook-substitutable-reasoning-d5f7734f79.pdf}
\caption[Substitutable reasoning]{\emph{Substitutable reasoning.} Two alternative engines use a domain's task and operation models while making different planning decisions. Model selection chooses a foundation model within an engine; controlled application and acceptance consume their own models and evidence. Type compatibility preserves the request form, not semantic quality or equivalent behavior.}
\label{handbook-substitutable-reasoning}
\end{figure}

% Evidence: BP-SUBSTITUTION-CONTRACTS | source: book/data/local-evidence/evidentiary-basis-11202026/integration/records/BP-SUBSTITUTION-CONTRACTS.md | kind: design-illustration
Consider an image-description task in the illustrated design. At the narrower substitution boundary, select a different foundation model while retaining the same task, permitted operations, and acceptance conditions. The surrounding workflow remains unchanged; only the component supplying the semantic judgment differs. At the broader boundary, a general-purpose agent may replace several specialized reasoning passes. It can choose which analyses to perform and in what order, but the design requires its proposed document changes to use the supported operation vocabulary and the resulting artifact to satisfy the applicable acceptance conditions. The first substitution changes a reasoner within a workflow; the second changes the workflow itself. Both retain the task, operation, and acceptance contracts.

\section{Engineering Models}\label{heading:engineering-models}

\autoref{handbook-substitution-models} identifies five models and their consumers. Task, operation, and acceptance models define the obligations that must remain stable across substitutions. The workflow model describes how reasoning is organized; the capability and selection model determines which foundation models are eligible for a task. The relationships among these models define the substitution boundaries.

{\footnotesize
\begin{longtable}[]{@{}
  >{\raggedright\arraybackslash}p{(\linewidth - 6\tabcolsep) * \real{0.2500}}
  >{\raggedright\arraybackslash}p{(\linewidth - 6\tabcolsep) * \real{0.2500}}
  >{\raggedright\arraybackslash}p{(\linewidth - 6\tabcolsep) * \real{0.2500}}
  >{\raggedright\arraybackslash}p{(\linewidth - 6\tabcolsep) * \real{0.2500}}@{}}
\caption[Engineering models for substitutable reasoning]{\emph{Engineering models for substitutable reasoning.} Task, operation, and acceptance models state the domain obligations while the reasoning workflow and selected foundation models may vary. Shared task and target identities connect these representations to their consumers.}\label{handbook-substitution-models}\tabularnewline
\toprule\noalign{}
\begin{minipage}[b]{\linewidth}\raggedright
Model
\end{minipage} & \begin{minipage}[b]{\linewidth}\raggedright
Engineering question
\end{minipage} & \begin{minipage}[b]{\linewidth}\raggedright
Representation
\end{minipage} & \begin{minipage}[b]{\linewidth}\raggedright
Relationships / consumers
\end{minipage} \\
\midrule\noalign{}
\endfirsthead
\toprule\noalign{}
\begin{minipage}[b]{\linewidth}\raggedright
Model
\end{minipage} & \begin{minipage}[b]{\linewidth}\raggedright
Engineering question
\end{minipage} & \begin{minipage}[b]{\linewidth}\raggedright
Representation
\end{minipage} & \begin{minipage}[b]{\linewidth}\raggedright
Relationships / consumers
\end{minipage} \\
\midrule\noalign{}
\endhead
\bottomrule\noalign{}
\endlastfoot
Task model & What judgment is required, and what makes its result acceptable? & Task identifiers, input/output shapes, semantic obligations, error consequences. & Joins the reasoning workflow to capability selection and task-specific evaluation. \\
Capability and selection model & Which reasoning resources are eligible for a task? & Model capabilities, task requirements, input traits, measured performance, selection policy. & Task identity joins requirements to candidate models; the selector derives a model choice. \\
Operation model & Which effects may the engine propose? & Typed operations, target references, arguments, preconditions, result meanings. & Targets refer to the domain artifact; the applier resolves and executes supported proposals. \\
Reasoning workflow model & How is work decomposed, ordered, and assigned? & A pass graph with dependency edges and task assignments, or an agent\textquotesingle s revisable plan. & Execution machinery consumes the workflow; its steps request tasks and propose operations. \\
Acceptance model & What evidence is required to accept the artifact? & Obligations, evaluation criteria, required evidence, blocking and review conditions. & Evaluators inspect the artifact; a gate or responsible reviewer consumes their findings. \\
\end{longtable}
}

\subsection*{Task, Operation, and Acceptance Contracts}

% Evidence: BP-SUBSTITUTION-CONTRACTS | source: book/data/local-evidence/evidentiary-basis-11202026/integration/records/BP-SUBSTITUTION-CONTRACTS.md | kind: implementation-and-design
The \textbf{operation contract} defines the changes a reasoning engine may request. DocAble's Office editing interface provides a closed vocabulary of operations, including setting alternative text, changing heading levels, and marking decorative content. Each operation identifies its target and required arguments. The application component resolves and executes supported operations under the applicable admissibility policy, returning application, refusal, or failure. \autoref{chapter:appendix-a-controlled-mutation} develops the corresponding boundary for PDF changesets, which use a different operation vocabulary.

The vocabulary must support legitimate work without becoming an unrestricted programming interface. A replacement engine can propose different operations or sequence them differently. To make this an enforcement boundary, the surrounding environment must also restrict mutation authority to the controlled application component. A typed interface alone does not prevent a reasoner with direct file access from bypassing it.

% Evidence: BP-SUBSTITUTION-CONTRACTS | source: book/data/local-evidence/evidentiary-basis-11202026/integration/records/BP-SUBSTITUTION-CONTRACTS.md | kind: implementation-and-design
The \textbf{task contract} defines a semantic judgment requested from a foundation model. DocAble identifies tasks such as image description, OCR transcription, reading-order inference, and heading adjudication, with declared output representations and capability requirements. Such definitions let a consumer request a task while a separate component handles model selection and invocation. They specify what the answer must mean, independently of which model supplies it.

The contract must also express what makes an answer acceptable. For alternative text, the obligation is not simply to produce a nonempty string: the description must convey the image's relevant information without inventing content. For decorative classification, an incorrect judgment may suppress meaningful content entirely. Such asymmetries influence model selection, validation, and the appropriate response to uncertainty.

The \textbf{acceptance contract} specifies the evidence required to accept the result. Deterministic checks can identify missing accessibility properties or illegal document structures. Content-preservation checks can detect selected forms of unintended loss. Semantic evaluators or human reviewers assess properties that structural checks cannot decide, such as whether an alternative description accurately conveys an image. The contract must also identify which findings block acceptance and which require engineering judgment.

\subsection*{Workflow, Selection, and Correspondence}

% Evidence: BP-SUBSTITUTION-CONTRACTS | source: book/data/local-evidence/evidentiary-basis-11202026/integration/records/BP-SUBSTITUTION-CONTRACTS.md | kind: implementation-and-design
The \textbf{reasoning workflow model} represents the process used to produce a result. DocAble organizes remediation as a graph of specialized passes. Some passes are algorithmic; others request structural judgments or descriptive content from foundation models. The graph represents dependencies among computations; execution machinery uses those dependencies to determine which work is ready as intermediate results are produced. A planning agent can instead maintain and revise a plan as it inspects the artifact. The workflow may change while task meanings and acceptance obligations remain fixed. A plan is guidance until an executor consumes it; an execution trace records the work actually performed.

Assign ownership to the obligations and to the observations used to evaluate them. The responsible engineer authorizes task meanings, permitted operations, and acceptance conditions; measured capability results describe behavior on an identified workload. A selection policy consumes those observations but does not turn them into universal quality guarantees. A selected model choice or a generated repair plan is derived from other representations. That describes its origin, not its authority.

Task identifiers join workflow steps to selection requirements and evaluators. Target identifiers join operation proposals to the current domain artifact. Keep these meanings aligned as schemas, models, and artifacts change: a task renamed without updating its evaluator, or a target resolved against stale structure, breaks correspondence even when the request remains well typed. Check that the executor follows required workflow dependencies and that the gate consumes the evidence the acceptance model requires.

The engine retains freedom to choose among admissible repairs, gather additional information, and revise their order where dependencies permit. Capability requirements and interface constraints restrict that freedom without prescribing every reasoning step.

\section{Realization}\label{heading:realization}

\subsection*{Connect the Models to Their Consumers}\label{heading:connect-the-models-to-their-consumers}

Construction begins with a bounded semantic task: its inputs, output representation, obligations, and evidence of success. The task definition must distinguish errors with different consequences, such as omitting useful alternative text and inventing content that was not present.

The application interface then establishes which effects a reasoning implementation may request. Supported operations have stable targets, explicit preconditions, and results distinguishing application, refusal, and failure. Restricting mutation authority to this interface prevents an alternative reasoner from bypassing the intended controls.

Acceptance requires separate evidence about the resulting artifact. Deterministic checks address mechanically decidable properties; preservation checks examine selected unintended changes; semantic evaluation addresses meaning. A gate must consume the required findings for them to affect acceptance. Only after these contracts are established can two reasoning engines be compared meaningfully: both must operate on comparable inputs through the same application and evaluation machinery.

\subsection*{Select Models for Fixed Tasks}\label{heading:select-models-for-fixed-tasks}

% Evidence: BP-SUBSTITUTION-CONTRACTS | source: book/data/local-evidence/evidentiary-basis-11202026/integration/records/BP-SUBSTITUTION-CONTRACTS.md | kind: implementation-and-design
% Evidence: BP-SUBSTITUTION-CONTRACTS | source: book/data/local-evidence/evidentiary-basis-11202026/integration/records/BP-SUBSTITUTION-CONTRACTS.md | kind: implementation-and-design
Model selection is the smaller substitution problem. A consumer requests a semantic task and supplies relevant input characteristics; the selector returns an eligible model, its provider, and the capabilities supporting the choice. It may expose a quality-tier rank for observability, but callers should not reconstruct selection or fallback policy from that rank.

Selection can account for task-specific capability requirements, input difficulty, provider availability, and cost. A capability requirement is a constraint: a model that cannot process the required input is not made suitable by being inexpensive or available. Among eligible models, a policy may prefer higher-quality alternatives, distribute requests among providers, or select a lower-cost model when its measured performance is adequate.

A model-selection policy should be evaluated against the errors that matter for each task. Average quality is not sufficient when different errors have substantially different consequences. For example, fabricating content in an alternative description and failing to recognize a decorative image are different failures, even if a benchmark combines them into one accuracy score.

Replacing the reasoning engine is a broader change: it may alter task decomposition, information gathering, operation ordering, and termination decisions.

\subsection*{Exercise an Alternative Engine}\label{heading:exercise-an-alternative-engine}

An alternative is to give a general-purpose agent the document, the available operations, and the acceptance criteria. The agent can inspect the document, run checks, decide which repairs to attempt, request supported operations, and inspect the resulting findings. Within the operation vocabulary's scope, it takes over planning and sequencing without requiring a new document-editing implementation.

% Evidence: E38 | source: book/data/local-evidence/evidentiary-basis-11202026/integration/records/E38.md | kind: dated-case-evidence
% Evidence: E38 | source: book/data/local-evidence/evidentiary-basis-11202026/integration/records/E38.md | kind: dated-case-evidence
% Evidence: E38 | source: book/data/local-evidence/evidentiary-basis-11202026/integration/records/E38.md | kind: dated-case-evidence

DocAble's substitution pilot gives a general-purpose agent access to existing document-checking, editing, and validation operations. The agent selects repairs and invokes those operations without changes to their deterministic implementations. The available operation vocabulary covers nearly all checker-classified fixable issues in the pilot corpus, although this measures interface expressiveness rather than successful remediation. A small subsequent comparison shows substantial differences in completion and remaining checker findings among foundation models. The pilot demonstrates reuse of the surrounding machinery, not equivalent performance across reasoning engines.

\subsection*{Evaluate the Result}\label{heading:evaluate-the-result}

A successful substitution must preserve the engineering obligations, not merely the interface.

% Evidence: E46 | source: book/data/local-evidence/evidentiary-basis-11202026/integration/records/E46.md | kind: illustration
Consider alternative text for a scientific figure. A reasoning engine may propose a valid operation that inserts a description, and the document checker may confirm that the image now has alternative text. Yet the description may repeat the figure's title without conveying its contents. The operation is valid and resolves the structural finding, but the accessibility requirement remains unsatisfied.

The same problem arises when deterministic heuristics stand in for a semantic judgment. In DocAble's alternative-text investigation, a repair strategy described as asking a language model whether an existing description is useful instead uses a chain of deterministic heuristics. An inadequate description, ``Heat Transfer 211,'' survives for a physics diagram. The decision procedure does not evaluate the semantic property promised by the contract. Deterministic methods can be appropriate substitutes, but only where evidence supports the property they must decide.

Separate evaluation makes reasoning strategies easier to compare, but does not guarantee independent errors. Structural and content-preservation checks supply evidence about the properties they inspect. Semantic judgments require evidence about meaning, which may come from another model or a human reviewer. An evaluator that shares the producer's model family or assumptions may confirm the same mistake.

A comparison should therefore use the same inputs and acceptance criteria for each candidate engine. Record the remaining violations, semantic quality, execution failures, latency, and cost. Where the evaluation cannot establish a property, retain that limitation rather than interpreting a passing checker as evidence of complete correctness.

% Evidence: E38 | source: book/data/local-evidence/evidentiary-basis-11202026/integration/records/E38.md | kind: dated-case-evidence
\autoref{chapter:5.2-inside-docables-software-factory} reports the experimental comparisons and their limitations.

For an end-to-end trial, introduce an unsupported operation, a stale target, a missing required evaluation result, and a plausible but inadequate description. Confirm the appropriate refusal or review, the artifact actually produced, and the evidence used for acceptance. A successful path must produce an acceptable artifact; a failed path must not acquire a success record merely because the engine reports completion.

\section{Adaptation and Limits}\label{heading:adaptation-and-limits}

Specialized reasoning workflows require construction and maintenance but can divide work into tasks that less expensive models perform adequately. General-purpose agents may require less specialized orchestration while consuming more inference resources. The economic choice depends on the quality required, the frequency of the task, and the costs of maintaining the surrounding machinery. As model capabilities change, an engineer may find it worthwhile to replace a specialized workflow or retain it for predictable, high-volume work.

The same architecture applies beyond document remediation. A planning agent can propose operations against a structured workflow; a classifier can produce typed decisions consumed by a controlled application; a diagnostic agent can recommend repairs whose effects are checked independently. In each case, the engineer determines which decisions may vary with the reasoning implementation and which obligations must remain stable.

The boundary need not be perfectly general. A small vocabulary supporting the operations a system actually requires is often more useful than a universal interface that restores unrestricted authority to the reasoner. New work may justify an extension, whose effects must be evaluated against the obligations already protected.
